\uuid{Q94V}
\titre{Pgcd}
\auteur{bodin ; croizat ; sacre}
\organisation{exo7}
\licence{CC BY-NC-SA 4.0}
\datecreate{2022-03-08}
\niveau{L1}
\module{Algèbre}
\chapitre{Arithmétique dans Z}
\sousChapitre{Pgcd, ppcm, algorithme d'Euclide}
\difficulte{2}
\theme{pgcd}
\exo7qcmid{hilisit.arithmetique.007}

\contenu{
\begin{question}
\qclassification{103.03}
    Grâce à l'application de l'algorithme d'Euclide, on obtient pour $a=630$ et $b=165$ :
    \begin{answers}
        \good{$\pgcd(a,b) = \pgcd(165,135)$}
        \good{$\pgcd(a,b) = \pgcd(135,30)$}
        \bad{$\pgcd(a,b) = \pgcd(30,0)$}
        \good{$\pgcd(a,b) = 15$}
    \end{answers}
    \begin{explanations}
     L'algorithme d'Euclide nous donne :\\
     $$ 630 = 165 \times 3 + 135 $$
     $$ 165 = 135 \times 1 + 30 $$
     $$ 135 = 30 \times 4 + \boxed{ \; 15 \; } $$
     $$ 30 = 15 \times 2 + 0 $$
     Ainsi on a :
     $$ \pgcd(a,b) = \pgcd(165,135) = \pgcd(135,30) = \pgcd(30,15) = \pgcd(15,0) = 15 $$
     En revanche, $ \pgcd(30,0) = 30 \neq \pgcd(a,b)$.
    \end{explanations}
\end{question}
}
